\scanpage{13}
\mspage{11}

un $S_n$ lisse sur $W_n$, de façon que ces choix se recollent (il suffit pour ceci S lisse et affine). Or, il résulte encore facilement des définitions, utilisant seulement que $S=S_0\longrightarrow S_n$ est une immersion fermée définie par un Idéal nilpotent, que le foncteur restriction des cristaux sur $S_n$ (relativement à une base quelconque -- ici on prendra $\unsure{R_n|W_n}$) vers les cristaux sur $S_0$, est une équivalence de catégories, donc un cristal sur S relativement à $\unsure{R_n|W_n}$ s'identifie à un cristal \inserted{sur $S_n$} relativement à $\unsure{R_n|W_n}$. Comme $S_n$ b est lisse sur $W_n$, il résulte de ce qui a été dit dans 1.8. que les cristaux en question s'identifient aux objets (de l'\unsure{espàce|espèce} $\underline F$ considérée) sur $S_n$, \uline{munis d'une stratification relativement à} $R_n=\operatorname{Spec}(W_n)$. Les foncteurs restrictions sur les cristaux s'expriment alors par les foncteurs restrictions correspondants pour objets stratifiés. En résumé: \uline{un cristal} \inserted{$\underline M$} \uline{sur} S \uline{s'identifie à un} \deleted{\unsure{pour|Pour}} \uline{système cohérent} $(\underline M_n)$ \uline{d'objets à stratification sur les différents} $S_n$ \uline{sur} $R_n=\operatorname{Spec}(W_n)$. Pour relier ceci à des objets qui soient plus dans la nature d'objets ``définis en caractéristique nulle'', supposons d'abord\unsure{s} que l'on puisse même relever S en un schéma X \uline{propre sur} $R=\operatorname{Spec}(W)$, hypothèse évidemment bien restrictive. On voit alors, utilisant les théorèmes de comparaison EGA III 4,5 que \uline{dans le cas de cristaux de modules dohérents sur} S, \uline{la catégorie desdits est canoniquement équivalente à la catégorie des Modules cohérents} $\underline M$ \uline{sur} X, \uline{munis d'une stratification relativement à} $R=\operatorname{Spec}(W)$. (NB Il se trouve donc, à posteriori, que cette \unsure{dernière|detnière} catégorie est essentiellement indépendante du relèvement choisi X de S). Considérant la fibre générique $X_K$ \deleted{sur} de X sur R, qui est un schéma propre et lisse sur un corps de caractéristique nulle, \unsure{on|oj} trouve donc un foncteur remarquable ``restriction'' ou ``localisation'', allant de la catégorie des cristaux de modules cohérents sur S, dans la catégorie des Modules cohérents \deleted{stratifiés} sur $X_K$, stratifiés relativement à $X_K$. Or (oubli de 1.8.) on voit \unsure{facilement} qu'un Module cohérent stratifié sur un préschéma localement de type fini sur un corps est nécessairement localement libre. De plus comme ici K est de caractéristique nulle \inserted{et $X_K$ lisse dessus}, on a signalé dans 1.8. que la notion de stratification s'explicite très simplement comme celle de ``connexion inté-
\margin{12}
